\section{递归地址化：派生概念由旧层读出序列生成}\label{sec:recursive-addressing}

本章并不是在第 \ref{sec:logic-expansion-chain} 节统一上位逻辑之外另行建立一套“地址逻辑”，而是把其中关于指称、局部性、\(\mathsf{NULL}\) 与观察者状态的抽象 forcing 骨架，具体实例化到递归地址化的对象层上。前一章已经通过 \(Fold_m\)、restriction 逆系统与多尺度一致性给出了稳定类型层的规范输出；本章则进一步回答：哪些旧层读出序列可以上升为地址，哪些地址只具有局部可读性，哪些地址虽然已经进入可接受引用域却仍必须判为 \(\mathsf{NULL}\)，以及一个地址何时能在 dyadic 证书链的逆极限中获得决定性的读出。因而，本章所完成的并不是“从类型到数值”的跨越，而是“从读出到对象”的跨越：先有可接受引用，后有前缀站点与局部截面，再有整体对象与读出决定性，最后才可能进入值层。

就本章内部结构而言，起首两小节把复合地址与递归语义写成旧层读出序列上的对象生成机制；“前缀站点与局部截面：\(\mathsf{NULL}\) 的内禀刻画”把第 \ref{sec:logic-expansion-chain} 节的局部对象—相容—粘接链具体化，从而使 \(\mathsf{NULL}\) 不再是特殊值，而成为地址层上的结构性失败；“观察者索引、信息状态与 forcing 语义”把地址对象重新放回状态 forcing 纤维；而 dyadic 证书链的逆极限寻址与读出决定性则给出对象何以在跨分辨率证书系统中进入决定包络的形式条件。于是，本章的总接口可以压缩为
\[
\text{旧层读出序列}
\Longrightarrow
\text{地址对象}
\Longrightarrow
\text{局部截面/\(\mathsf{NULL}\) 判定}
\Longrightarrow
\text{逆极限决定性}.
\]
这也意味着：本章是全文的对象层实例章，它回答的是“何物已经成为对象”，而不是提前回答“该对象取何值”。


本节在标准可测动力系统 \((X,\mathcal{B},\mu,T)\) 与有限时间窗读出口径下处理一个结构性问题：在给定一族可测概念之后，如何由其有限时间窗读出递归生成更细的概念族，以及该过程是否会在可测结构层面引入外加对象。核心结论是：递归地址化不扩张既有可见 \(\sigma\)-代数；递归生成的柱事件族仍落在旧层可见 \(\sigma\)-代数 \(\mathcal{G}^{(L)}\) 内（命题 \ref{prop:recursive-addressing-refinement}），从而新层仅在 \(\mathcal{G}^{(L)}\) 内重新选取并组织可见事件（推论 \ref{cor:recursive-addressing-visible-sigma-nonexpansion}，因而 \(\mathcal{G}^{(L+1)}\subseteq \mathcal{G}^{(L)}\)）。在符号柱集口径下，前缀细化对应前缀超度量球的缩小，见第 \ref{sec:spg} 节段落 \ref{par:spg-prefix-ultrametric}。与第 \ref{sec:logic-expansion-chain} 节的抽象扩张链相对应，本节把地址、局部截面、\(\mathsf{NULL}\) 与观察者索引 forcing 解释为该上位语义在扫描—投影读出模型中的具体实现。

\subsection{地址的复合生成}

设在层级 $L$ 已得到一族概念 $\mathcal{C}^{(L)}=\{C^{(L)}_1,\dots,C^{(L)}_{r_L}\}$，其中每个 $C^{(L)}_i\in\mathcal{B}$ 可测。定义派生的二元向量过程
\[
\mathbf{b}^{(L)}_t(x):=\bigl(\ind_{C^{(L)}_1}(T^t x),\dots,\ind_{C^{(L)}_{r_L}}(T^t x)\bigr)\in\{0,1\}^{r_L}.
\]
对任意长度 $\ell$ 与任意地址
\[
a=\bigl(a_0,\dots,a_{\ell-1}\bigr)\in\bigl(\{0,1\}^{r_L}\bigr)^\ell,
\]
定义其对应的派生柱事件
\[
C^{(L+1)}_{a,t}:=\{x\in X:\mathbf{b}^{(L)}_t(x)=a_0,\dots,\mathbf{b}^{(L)}_{t+\ell-1}(x)=a_{\ell-1}\}\in\mathcal{B}.
\]
令 $\mathcal{C}^{(L+1)}$ 为选定的一族此类柱事件（例如取一组有限地址集合 $A^{(L+1)}$ 对应的 $C^{(L+1)}_{a,0}$）。由此得到闭合递归：\textbf{派生概念是旧层读出序列的柱事件}，从而地址仅来自可观测数据而非外加原语。

\begin{proposition}[$\sigma$-代数细化：递归地址化不引入外加结构]\label{prop:recursive-addressing-refinement}
令 $\mathcal{G}^{(L)}$ 为由旧层概念在有限时间窗内的平移生成的 $\sigma$-代数：
\[
\mathcal{G}^{(L)}:=\sigma\bigl(T^{-t}C_i^{(L)}:\ 1\le i\le r_L,\ t\in\ZZ\bigr).
\]
则对任意地址 $a$ 与任意 $t$，派生柱事件 $C^{(L+1)}_{a,t}$ 属于 $\mathcal{G}^{(L)}$。因此递归地址化仅执行“对既有可观测事件作柱集闭包/细化”，而不引入外加可测集。
\leanverified{paper\_recursive\_addressing\_refinement}
\end{proposition}

\begin{proof}
对任意固定的 \(i\) 与 \(t\)，函数 \(x\mapsto \ind_{C_i^{(L)}}(T^t x)=\ind_{T^{-t}C_i^{(L)}}(x)\) 是 \(\mathcal{G}^{(L)}\)-可测的。
因此对任意给定地址 \(a=(a_0,\dots,a_{\ell-1})\)，事件
\[
\{x:\mathbf b^{(L)}_t(x)=a_0,\dots,\mathbf b^{(L)}_{t+\ell-1}(x)=a_{\ell-1}\}
\]
可写为有限个集合 \(T^{-s}C_i^{(L)}\) 及其补集的有限交（在 \(s=t,\dots,t+\ell-1\) 上并对 \(i=1,\dots,r_L\) 展开），故属于由这些集合生成的 \(\sigma\)-代数，即属于 \(\mathcal{G}^{(L)}\)。证毕。
\end{proof}

\begin{corollary}[可见 \texorpdfstring{$\sigma$}{sigma}-代数不扩张：递归柱事件仅在内层重组织]\label{cor:recursive-addressing-visible-sigma-nonexpansion}
沿用命题 \ref{prop:recursive-addressing-refinement} 的记号。
设新层概念族 \(\mathcal{C}^{(L+1)}\) 由若干派生柱事件 \(C^{(L+1)}_{a,t}\) 选取而成，并令
\[
\mathcal{G}^{(L+1)}:=\sigma\bigl(T^{-t}C:\ C\in \mathcal{C}^{(L+1)},\ t\in\ZZ\bigr).
\]
则有
\[
\boxed{\ \mathcal{G}^{(L+1)}\subseteq \mathcal{G}^{(L)}\ }.
\]
因此，递归生成的柱事件族不扩张既有可见 \(\sigma\)-代数，而仅在 \(\mathcal{G}^{(L)}\) 内重新选取并组织可见事件。
\leanverified{paper\_recursive\_addressing\_visible\_sigma\_nonexpansion}
\end{corollary}

\begin{proof}
由命题 \ref{prop:recursive-addressing-refinement}，对任意地址 \(a\) 与任意 \(t\in\ZZ\)，派生柱事件 \(C^{(L+1)}_{a,t}\) 都属于 \(\mathcal{G}^{(L)}\)。
设 \(C\in\mathcal{C}^{(L+1)}\)，按构造存在 \(a,t_0\) 使 \(C=C^{(L+1)}_{a,t_0}\)。
于是对任意 \(s\in\ZZ\)，有 \(T^{-s}C=T^{-s}C^{(L+1)}_{a,t_0}=C^{(L+1)}_{a,t_0+s}\in\mathcal{G}^{(L)}\)。
因此 \(\mathcal{G}^{(L+1)}\) 的生成元均属于 \(\mathcal{G}^{(L)}\)，从而 \(\mathcal{G}^{(L+1)}\subseteq \mathcal{G}^{(L)}\)。证毕。
\end{proof}

\subsection{空值到投影的递归语义}

与推论 \ref{cor:recursive-addressing-visible-sigma-nonexpansion} 的“可见 \texorpdfstring{$\sigma$}{sigma}-代数不扩张”结论相一致，递归地址化在语义上可表述为“地址先于取值”的\emph{可读出前缀}机制。对任意候选地址 \(a\in A^{(L+1)}\)，在地址对象未给定之前不对其派生概念赋值：读出为偏函数，在该输入处不定义，其失败不得被误读为数值 \(0\)，并按本文约定以内禀空值 \(\mathsf{NULL}\) 表示“不可寻址/不定义”。一旦地址 \(a\) 被选定，则概念值由二元投影 \(x\mapsto \ind_{C_a^{(L+1)}}(x)\) 给出。于是层级递归满足：
\begin{itemize}
  \item “无地址 $\Rightarrow$ 空值（不定义）”；
  \item “有地址 $\Rightarrow$ 二元投影（$0/1$）”；
  \item “地址 $\Rightarrow$ 由旧层读出生成”。
\end{itemize}
在类型语义上，这意味着：若地址对象未给定（或不属于协议允许的 \(\mathcal{A}_{\mathrm{adm}}\)），则不存在可在对象层引用的评价函数 \(\ind_{C_a}\)（等价地，\(1_{C_a}\) 不能在协议的可见域中被构造成一个总定义读出）；
此时读出偏函数在该输入处不定义，其语义后果以内禀空值 \(\mathsf{NULL}\) 表示（并将在 \S\ref{subsec:prefix-site-null} 中被等价内生化为 \(\mathscr{F}(a)=\varnothing\) 的局部截面障碍）。
本节所述空值属于\emph{语义空值}（地址缺失，读出不定义）；与之正交的\emph{协议空值}（类型系统锁定 unitary slice）与\emph{碰撞空值}（侧信息预算不足导致不可唯一重建）见 \S\ref{subsec:xi-pw-type-safety-null} 及命题 \ref{prop:xi-null-three-way}。在更统一的口径下，可把“完备地址”打包为多轴前缀（定义 \ref{def:multi_axis_address}），其中碰撞空值对应残差轴不足，其最小侧信息预算由定理 \ref{thm:address-defect-law} 给出。
该语义可与布尔代数闭包兼容，也可在算子代数层面以投影算子及其乘积实现（见附录 \ref{app:op-algebra}）。

\subsection{前缀站点与局部截面：\texorpdfstring{$\mathsf{NULL}$}{NULL} 的内禀刻画}\label{subsec:prefix-site-null}
本小节给出一个与“地址先于取值”严格兼容的形式化：把地址空间视为由前缀偏序生成的站点（cylinder site），并把“协议可接受且单值可核验的读出/证书”组织为其上的（预）层对象。其用途不在于引入新的测度论结构，而在于把空值（\(\mathsf{NULL}\)）提升为一个清晰的存在性障碍：\emph{在给定地址处不存在协议允许的局部截面}，从而避免把 “不定义” 误读为数值 \(0\) 的歧义。

\begin{definition}[前缀站点与柱覆盖]\label{def:prefix-site}
为简明起见仅以二元字母表 \(\Sigma=\{0,1\}\) 书写（对一般有限字母表 \(\Sigma\) 完全同型）。令 \(\mathcal{A}:=\Sigma^{<\infty}\) 为全体有限地址，并以前缀关系 \(a\preceq b\) 赋偏序。
令 \(\mathcal{A}_{\mathrm{adm}}\subseteq\mathcal{A}\) 为协议所允许的地址集合（在本节的递归构造中，其原型可取为各层选定的有限地址族 \(A^{(L)}\) 及其在 refinement 下的闭包）。
对任意 \(a\in\mathcal{A}_{\mathrm{adm}}\)，令柱事件 \(C_a\subseteq X\) 如第 \ref{sec:spg} 节定义（亦即 \(C_a=\omega_\infty^{-1}([a])\)，其中 \([a]\subseteq\Omega_\infty\) 为符号柱集）。
\par
称有限族 \(\{b_i\succeq a\}_{i=1}^n\subseteq\mathcal{A}_{\mathrm{adm}}\) 为 \(a\) 的一个\textbf{柱覆盖}，若 \(C_a=\bigcup_{i=1}^{n} C_{b_i}\)。
该覆盖关系与段落 \ref{par:spg-prefix-ultrametric} 的前缀拓扑一致：\(\Omega_\infty=\varprojlim \Omega_k\) 的 clopen 基由柱集 \([a]\) 给出，因而“前缀细化”与“覆盖分解”为同一对象的两种等价描述。
\end{definition}

\begin{definition}[协议可接受读出的（预）层]\label{def:prefix-sheaf-readout}
固定一套“可见域 + 检查器”的读出协议，其抽象口径见定义 \ref{def:xi-visible-read-null}（在此仅使用其“读出可为部分映射、并以 \(\mathsf{NULL}\) 表示不定义”的结构）。
对每个 \(a\in\mathcal{A}_{\mathrm{adm}}\) 定义集合
\[
\mathscr{F}(a):=\{\text{在柱事件 }C_a\text{ 上协议允许、且在类型意义下为单值可核验的读出/证书对象}\},
\]
并对 refinement \(b\succeq a\) 定义限制映射
\(\rho_{a\leftarrow b}:\mathscr{F}(a)\to\mathscr{F}(b)\)
为把读出/证书限制到子柱 \(C_b\)（等价地：在事件层面对更深前缀条件化）。
\end{definition}

\begin{proposition}[局部截面存在性与 \texorpdfstring{$\mathsf{NULL}$}{NULL}]\label{prop:null-as-local-section-obstruction}
在定义 \ref{def:prefix-site}--\ref{def:prefix-sheaf-readout} 的口径下，“地址先于取值”的语义可等价表述为：
\[
\text{地址 }a\text{ 的读出在类型层非 }\mathsf{NULL}
\quad\Longleftrightarrow\quad
\mathscr{F}(a)\neq\varnothing.
\]
并且在递归地址化的闭包口径下（命题 \ref{prop:recursive-addressing-refinement}），\(\mathscr{F}\) 对任意有限柱覆盖满足粘接唯一性：若 \(\{b_i\succeq a\}\) 覆盖 \(a\)，且一族局部对象 \(s_i\in\mathscr{F}(b_i)\) 在重叠处相容（即在任意交 \(C_{b_i}\cap C_{b_j}\) 上的限制一致），则存在唯一 \(s\in\mathscr{F}(a)\) 使得 \(s|_{C_{b_i}}=s_i\)。
\leanverified{paper\_null\_as\_local\_section\_obstruction}
\end{proposition}

\begin{proof}
第一条等价仅是将 “不定义（\(\mathsf{NULL}\)）” 从数值层改写为对象层的存在性断言：读出/证书之所以可核验，要求其在给定地址处确有协议允许的局部对象可供引用；当且仅当该局部对象存在时，读出在类型层被视为非 \(\mathsf{NULL}\)。
\par
对粘接唯一性：柱覆盖是有限并集分解，且柱事件在 refinement 下封闭（由定义）；递归地址化又保证所有派生对象仍落在既有可观测柱事件所生成的 \(\sigma\)-代数内（命题 \ref{prop:recursive-addressing-refinement}）。
因此若局部对象在重叠处一致，则它们在集合层的拼接由并集给出且唯一；协议可接受性在限制映射下保持，故粘接得到的对象仍属于 \(\mathscr{F}(a)\) 并满足唯一性。
\end{proof}

\begin{corollary}[\texorpdfstring{$\mathsf{NULL}$}{NULL} 的三分解（局部截面障碍版）]\label{cor:null-trichotomy-local-section}
在上述站点/（预）层表述下，\(\mathsf{NULL}\) 的三类来源（定义 \ref{def:xi-visible-read-null} 与命题 \ref{prop:xi-null-three-way}）可被统一为三类互不混淆的局部截面障碍：
\begin{enumerate}
  \item \textbf{语义空值：} \(a\notin\mathcal{A}_{\mathrm{adm}}\)，即协议未提供该地址对象；此时在站点上无对象可供取截面，空值由“缺乏指称”直接给出。
  \item \textbf{协议空值：} \(a\in\mathcal{A}_{\mathrm{adm}}\) 但可见域/类型系统拒绝在 \(C_a\) 上建立读出（典型地：可见域被锁定为某一切片而拒绝额外坐标轴），从而 \(\mathscr{F}(a)=\varnothing\)。
  \item \textbf{碰撞空值：} \(a\in\mathcal{A}_{\mathrm{adm}}\) 且通过检查器，但在“单值可核验读出”的类型约定下，若底层读出存在多对一压缩而未外置使扩展读出单射化的残差轴，则在该类型中仍等价表现为 \(\mathscr{F}(a)=\varnothing\)；其最小侧信息预算由命题 \ref{prop:xi-null-three-way} 与定理 \ref{thm:address-defect-law} 给出。
\end{enumerate}
因此，“\(\mathsf{NULL}\)”并非字符串的数值性质，而是站点对象与（预）层纤维在协议意义下的结构性空缺：合法性由 \(\mathscr{F}(a)\) 的非空性判定。
\leanverified{paper\_null\_trichotomy\_local\_section}
\end{corollary}

\paragraph{粘合障碍的上同调口径：\texorpdfstring{$H^2$}{H2} 见证}
命题 \ref{prop:null-as-local-section-obstruction} 将 \(\mathsf{NULL}\) 内禀化为“局部截面不存在”的对象论断言。
在本文后续的多轴前缀语义（定义 \ref{def:multi_axis_address}）与协议闭包中，局部证书对象往往还携带一组\emph{可调但必须可审计}的相容性标签（例如纤维自同构/gauge 选择的等价类；可取其阿贝尔化作为系数，见定义 \ref{def:fold-gauge-group}）。
此时，“在重叠处相容”不仅是集合意义的相等，还可能包含一个可复核的二阶一致性约束；其失败可用 Čech 上同调的二阶障碍类作证书化刻画。

\begin{definition}[粘合障碍类（Čech \texorpdfstring{$H^2$}{H2} 口径）]\label{def:prefix-site-h2-gluing-obstruction}
设 \(A\) 为一个系数阿贝尔群。
考虑上述前缀站点 \(\mathcal{A}_{\mathrm{adm}}\) 及其柱覆盖。
假设在某一协议口径下，局部证书对象的相容性可被组织为一个带 \(A\)-纤维的“相容层”（例如 \(A\)-torsor/带 gauge 的局部对象），使得对任意覆盖 \(\{b_i\succeq a\}\) 与任意局部对象族 \((s_i\in\mathscr{F}(b_i))\)，其在二重重叠上的差异可由一族 \(A\)-值转移数据 \(g_{ij}\) 记录，并在三重重叠上诱导一个 Čech \(2\)-余圈
\[
\alpha_{ijk}:=g_{ij}+g_{jk}+g_{ki}\in A,
\qquad
\alpha\in Z^2(\{b_i\},A).
\]
称其同调类 \([\alpha]\in H^2(\{b_i\},A)\)（并在加细覆盖下取极限）为地址 \(a\) 处的\textbf{粘合障碍类}：\([\alpha]=0\) 表示可通过允许的局部 gauge 调整把 \((s_i)\) 改写为在重叠处严格相容的家族，从而满足可粘接性；\([\alpha]\neq 0\) 则表示存在结构性扭结阻止其全球化。
\end{definition}

\begin{proposition}[\texorpdfstring{$H^2$}{H2} 粘合障碍 \(\Rightarrow\) \texorpdfstring{$\mathsf{NULL}$}{NULL}]\label{prop:null-as-h2-obstruction}
在定义 \ref{def:prefix-site-h2-gluing-obstruction} 的设定下，
若某一可核对地址 \(a\in\mathcal{A}_{\mathrm{adm}}\) 的粘合障碍类 \([\alpha]\neq 0\)，
则该地址处不存在协议允许的全局局部截面（证书对象），从而
\(
\mathscr{F}(a)=\varnothing
\)
并且读出在类型层必为 \(\mathsf{NULL}\)（命题 \ref{prop:null-as-local-section-obstruction}）。
\leanverified{paper\_recursive\_addressing\_null\_as\_h2\_obstruction}
\end{proposition}

\begin{proof}
反设 \(\mathscr{F}(a)\neq\varnothing\)，取任一全局对象 \(s\in\mathscr{F}(a)\) 并限制到覆盖 \(\{b_i\}\) 上得 \(s_i:=s|_{C_{b_i}}\in\mathscr{F}(b_i)\)。
则在任意重叠处这些限制严格相容；按定义 \ref{def:prefix-site-h2-gluing-obstruction}，其诱导的转移数据可取为平凡（从而 \(\alpha\) 为零余圈、\([\alpha]=0\)），与 \([\alpha]\neq 0\) 矛盾。
故 \(\mathscr{F}(a)=\varnothing\)，进而读出为 \(\mathsf{NULL}\)。
证毕。
\end{proof}
\leanverified{paper\_null\_as\_h2\_obstruction}

\paragraph{从 \(H^2\) 障碍到带状 gerbe：对象层级的提升}
命题 \ref{prop:null-as-h2-obstruction} 把 \([\alpha]\neq 0\) 解释为“全局证书对象不存在”的见证。
为了在结构上保留“局部存在、全局失效”的信息，单纯的集合值（预）层口径并不足够：当粘合在三重重叠上出现二阶结合缺陷时，合适的语义对象是以 \(A\) 为 band 的 \(2\)-层级对象。

\begin{theorem}[Čech-\texorpdfstring{$H^2$}{H2} 障碍的 gerbe 语义]\label{thm:prefix-site-cech-null-gerbe}
沿用定义 \ref{def:prefix-site-h2-gluing-obstruction} 的设定。
固定 \(a\in\mathcal{A}_{\mathrm{adm}}\) 与其一组柱覆盖 \(\{b_i\succeq a\}\)，并记二重与三重交为
\(
C_{b_{ij}}:=C_{b_i}\cap C_{b_j}
\)
与
\(
C_{b_{ijk}}:=C_{b_i}\cap C_{b_j}\cap C_{b_k}
\)。
设 \(\alpha\in Z^2(\{b_i\},A)\) 为该覆盖上由局部识别数据诱导的 Čech \(2\)-余圈，并记其同调类为 \([\alpha]\)（在加细覆盖下取极限）。
则存在一个由 \((\{b_i\},\alpha)\) 函子性构造的纤维化群胚 \(\mathcal{G}_\alpha\to \mathcal{A}_{\mathrm{adm}}\)，满足：
\begin{enumerate}
  \item \(\mathcal{G}_\alpha\) 为以 \(A\) 为 band 的带状 gerbe：在足够细覆盖下局部非空且局部连通，并且任意对象的自同构群自然同构于 \(A\)。
  \item \(\mathcal{G}_\alpha\) 的中性化（存在全局对象）与“地址 \(a\) 处存在协议允许的全局证书对象”严格等价。
  \item \(\mathcal{G}_\alpha\) 的等价类仅由 \([\alpha]\) 决定：更精确地，\(\mathcal{G}_\alpha\) 的带状类与 \([\alpha]\) 在 \(H^2(\{b_i\},A)\) 中一致。
\end{enumerate}
\end{theorem}

\begin{proof}
以覆盖索引集为对象，令 \(\mathcal{G}_\alpha(b_i)\) 的对象为“定义在 \(C_{b_i}\) 上的局部证书对象”，态射为协议允许的局部 gauge 变换。
在二重交 \(C_{b_{ij}}\) 上以转移数据 \(g_{ij}\) 识别局部对象，并要求在三重交 \(C_{b_{ijk}}\) 上的两种复合 \(g_{ij}g_{jk}\) 与 \(g_{ik}\) 的差异由 \(\alpha_{ijk}\) 给出（以 \(A\) 的作用实现结合缺陷）。
局部对象的存在性给出局部非空性；在假设“局部证书对象带 \(A\)-纤维的相容结构”下，任意两局部对象在足够细覆盖下局部同构，从而得局部连通性；而自同构的局部群由 \(A\) 决定并在加细下稳定，故得到 \(A\)-带状 gerbe。
\par
若存在全局证书对象，则其限制与恒等粘合数据给出中性对象；反之，中性对象给出覆盖上的 descent datum，并且中性化意味着可在允许的 gauge 调整下将结合缺陷严格消去，从而得到全局对象。
最后，更换局部对象与识别只会把 \(\alpha\) 改变为 Čech 余边界，因此 \(\mathcal{G}_\alpha\) 的等价类仅依赖于 \([\alpha]\)。
证毕。
\end{proof}
\leanverified{paper\_recursive\_addressing\_prefix\_site\_cech\_null\_gerbe}
\leanverified{paper\_prefix\_site\_cech\_null\_gerbe}

\paragraph{障碍的最小可见化：通用系数商}
当 \([\alpha]\neq 0\) 时，中性化失败把读出推向 \(\mathsf{NULL}\)（命题 \ref{prop:null-as-h2-obstruction}）。
另一方面，若允许把系数群作最小化商化，则可在不改变覆盖与局部数据的前提下把二阶缺陷严格消去，从而将 gerbe 归约到中性情形。

\begin{proposition}[使障碍点态消失的最小系数商]\label{prop:prefix-site-min-visible-quotient}
在定理 \ref{thm:prefix-site-cech-null-gerbe} 的设定下，令
\[
H_\alpha:=\bigl\langle \alpha_{ijk}\ \bigm|\ C_{b_{ijk}}\neq\varnothing \bigr\rangle\le A,
\qquad
A_{\mathrm{vis}}:=A/H_\alpha,
\]
并记自然投影为 \(\pi:A\twoheadrightarrow A_{\mathrm{vis}}\)。
则：
\begin{enumerate}
  \item \(\alpha\) 推送到 \(A_{\mathrm{vis}}\) 后点态为零：其像 \(\bar\alpha\in Z^2(\{b_i\},A_{\mathrm{vis}})\) 满足 \(\bar\alpha_{ijk}=0\) 对一切 \(C_{b_{ijk}}\neq\varnothing\)。
  因而对应的 gerbe \(\mathcal{G}_{\bar\alpha}\) 必为中性 gerbe。
  \item 若 \(\varphi:A\to B\) 为任意阿贝尔群同态，使得 \(\varphi(\alpha_{ijk})=0\) 对所有三重交成立（等价地，\(\alpha\) 在系数 \(B\) 下点态消失），则 \(\varphi\) 唯一因子分解为
  \[
  A\twoheadrightarrow A_{\mathrm{vis}}\xrightarrow{\ \exists!\ } B.
  \]
\end{enumerate}
\end{proposition}

\begin{proof}
(1) 由 \(H_\alpha\) 的生成方式，\(\pi(\alpha_{ijk})=0\) 对所有 \(C_{b_{ijk}}\neq\varnothing\) 成立，故 \(\bar\alpha=0\)。
在定理 \ref{thm:prefix-site-cech-null-gerbe} 的构造中，\(\bar\alpha=0\) 对应平凡结合子，从而 gerbe 中性。
\par
(2) 若 \(\varphi(\alpha_{ijk})=0\) 对一切三重交成立，则 \(H_\alpha\subseteq\ker\varphi\)，因此 \(\varphi\) 通过商 \(A/H_\alpha\) 因子分解且唯一。
证毕。
\end{proof}
\leanverified{paper\_recursive\_addressing\_prefix\_site\_min\_visible\_quotient}
\leanverified{paper\_prefix\_site\_min\_visible\_quotient}

\paragraph{保留分辨率的代价：外置寄存器的严格下界}
命题 \ref{prop:prefix-site-min-visible-quotient} 表明：在只允许改变系数群而不引入外置记录轴的制度下，\([\alpha]\) 的可见化最小化为 \(A_{\mathrm{vis}}\)。
若目标是恢复原始 \(A\) 的分辨率，则必须显式外置承载 \(H_\alpha\) 的信息量。

\begin{theorem}[恢复 \(A\)-分辨率所需的最小寄存器容量]\label{thm:prefix-site-register-lower-bound}
设 \(A\) 为有限阿贝尔群，并沿用命题 \ref{prop:prefix-site-min-visible-quotient} 的记号。
若仅允许额外引入一个有限寄存器集合 \(R\)，并要求存在集合嵌入
\[
A\hookrightarrow A_{\mathrm{vis}}\times R,
\]
则必有
\[
\card{R}\ \ge\ \card{H_\alpha}
\qquad\Longleftrightarrow\qquad
\log_2\card{R}\ \ge\ \log_2\card{H_\alpha}.
\]
并且该下界可达：存在 \(R\) 使得等号成立。
\leanverified{paper\_prefix\_site\_register\_lower\_bound}
\end{theorem}

\begin{proof}
对任意 \(\bar a\in A_{\mathrm{vis}}\)，投影 \(\pi:A\to A_{\mathrm{vis}}\) 的纤维 \(\pi^{-1}(\bar a)\) 为一个 \(H_\alpha\)-陪集，故其基数恒为 \(\card{H_\alpha}\)。
若存在嵌入 \(A\hookrightarrow A_{\mathrm{vis}}\times R\)，则其在每个纤维上的限制给出
\(
\pi^{-1}(\bar a)\hookrightarrow \{\bar a\}\times R\cong R
\)，
从而 \(\card{R}\ge \card{H_\alpha}\)。
\par
可达性：取 \(R:=H_\alpha\)。
任选一个集合截面 \(s:A_{\mathrm{vis}}\to A\)（不必为群同态），则每个 \(a\in A\) 可唯一写为 \(a=s(\pi(a))+h(a)\) 且 \(h(a)\in H_\alpha\)。
令 \(\iota(a):=(\pi(a),h(a))\in A_{\mathrm{vis}}\times H_\alpha\)，则 \(\iota\) 为嵌入且 \(\card{R}=\card{H_\alpha}\)。
证毕。
\end{proof}

\paragraph{代数化的读出：扭曲群胚与投影表象}
二阶结合缺陷亦可被转写为覆盖神经上的扭曲乘法，从而把“局部证书 + 识别数据”的读出链条解释为投影表象而非严格表象。

\begin{proposition}[Čech 群胚的 \(A\)-扭曲与投影表象]\label{prop:prefix-site-projective-rep-twist}
令 \(\mathfrak N(\{b_i\})\) 为覆盖的神经，并令 \(\Pi\) 为其 \(1\)-骨架生成的 Čech 群胚（对象为指标 \(i\)，态射由二重交 \(C_{b_{ij}}\neq\varnothing\) 生成）。
则 \(\alpha\) 诱导：
\begin{enumerate}
  \item 一个以 \(A\) 为中心的 \(\Pi\) 的 \(2\)-上扩张（等价地，\(\Pi\) 上的 \(A\)-加权结合子），从而定义扭曲群胚代数 \(\CC_\alpha[\Pi]\)（其乘法由 \(\alpha\) 扭曲）。
  \item 覆盖上的任一 descent 数据（局部对象与二重交识别）自然给出 \(\Pi\) 的一个投影表象，其乘子（multiplier）由 \([\alpha]\) 决定；并且该投影表象可严格化为真实表象当且仅当 \([\alpha]=0\)。
\end{enumerate}
\end{proposition}

\begin{proof}
在 \(\Pi\) 上选择生成态射 \(i\to j\) 表示二重交上的识别。
三重交上两种复合 \(i\to j\to k\) 与 \(i\to k\) 的差异由 \(A\)-值元素 \(\alpha_{ijk}\) 记录，这正是 \(2\)-范畴意义下结合子的全部数据。
将该结合子视为乘法扭曲即得到 \(\CC_\alpha[\Pi]\)。
\par
对投影表象：任一把局部对象与识别送入线性化实现的选择，在复合路径上得到的算子只在标量乘子（或一般的 \(A\) 作用）意义下相等，乘子由 \(\alpha_{ijk}\) 给定并在余边界改变下保持同调类不变。
当且仅当 \([\alpha]=0\) 时，存在余边界修正消除该乘子，从而 strictification 为真实表象。
证毕。
\leanverified{paper\_prefix\_site\_projective\_rep\_twist}
\leanverified{paper\_recursive\_addressing\_prefix\_site\_projective\_rep\_twist}
\end{proof}

\paragraph{对偶下降与谱可定义性：角色通道的可见子群}
在角色分解口径下，只有能湮灭 \(H_\alpha\) 的角色才可在可审计意义下给出无歧义的全局模式。

\begin{corollary}[可见角色与湮灭子群]\label{cor:prefix-site-visible-characters-annihilator}
设 \(A\) 为有限阿贝尔群，并记 Pontryagin 对偶 \(\widehat A:=\Hom(A,\TT)\)。
令 \(H_\alpha\le A\) 与 \(A_{\mathrm{vis}}=A/H_\alpha\) 如命题 \ref{prop:prefix-site-min-visible-quotient}。
则有自然同构
\[
\widehat{A_{\mathrm{vis}}}\ \cong\ H_\alpha^\perp
:=\{\chi\in\widehat A:\ \chi(h)=1\ \forall h\in H_\alpha\}.
\]
并且 \(\chi\in H_\alpha^\perp\) 当且仅当 \(\chi(\alpha_{ijk})=1\) 对所有 \(C_{b_{ijk}}\neq\varnothing\) 成立。
\end{corollary}

\begin{proof}
商同态 \(A\twoheadrightarrow A/H_\alpha\) 的对偶同态为嵌入 \(\widehat{A/H_\alpha}\hookrightarrow \widehat A\)，其像恰为湮灭子群 \(H_\alpha^\perp\)。
又因 \(H_\alpha\) 由 \(\{\alpha_{ijk}\}\) 生成，\(\chi\in H_\alpha^\perp\) 等价于 \(\chi(\alpha_{ijk})=1\) 对所有生成元成立。
证毕。
\end{proof}
\leanverified{paper\_recursive\_addressing\_prefix\_site\_visible\_characters\_annihilator}

\begin{theorem}[谱化输出的可定义性判别：只能在 \texorpdfstring{$A_{\mathrm{vis}}$}{Avis} 上无歧义成立]\label{thm:prefix-site-spectral-definability-avis}
考虑任一“谱化”步骤，其输入为命题 \ref{prop:prefix-site-projective-rep-twist} 的 descent 数据（等价地，扭曲群胚代数 \(\CC_\alpha[\Pi]\)），输出为可审计的谱对象（例如：转移算子族、迹/行列式、\(\zeta\)-函数型不变量或基于角色分解的误差控制）。
假设该步骤满足最弱的自然性：对允许的局部 gauge 调整（即改变 \(\alpha\) 为同调余边界）最终输出不变。
则：
\begin{enumerate}
  \item 若 \(\chi\in\widehat A\) 不属于 \(H_\alpha^\perp\)，则任何以 \(\chi\) 作为模式的“扭曲谱对象”都不具备 gauge 不变性，从而无法下降为全局可审计对象。
  \item 若 \(\chi\in H_\alpha^\perp\cong\widehat{A_{\mathrm{vis}}}\)，则上述歧义消失；因此一切基于角色分解的谱命题在形式上只能理解为关于 \(A_{\mathrm{vis}}\) 的命题，而非关于原始 \(A\) 的命题。
\end{enumerate}
\end{theorem}

\begin{proof}
由命题 \ref{prop:prefix-site-projective-rep-twist}，沿不同复合路径得到的线性化实现相差乘子 \(\chi(\alpha_{ijk})\)。
若存在某个三重交使 \(\chi(\alpha_{ijk})\neq 1\)，则该乘子在 gauge 调整下不能被消去，从而违背“最终输出 gauge 不变”的可审计性要求，故全局对象不可定义。
反之，当 \(\chi\in H_\alpha^\perp\) 时，对所有三重交均有 \(\chi(\alpha_{ijk})=1\)，投影乘子消失，谱化输出可下降为无歧义的全局对象；再由推论 \ref{cor:prefix-site-visible-characters-annihilator} 得到其必经由 \(\widehat{A_{\mathrm{vis}}}\)。
证毕。
\end{proof}
\leanverified{paper\_recursive\_addressing\_prefix\_site\_spectral\_definability\_avis}

\subsection{观察者索引、信息状态与 forcing 语义}\label{subsec:observer-indexed-forcing}
前述"地址先于取值""\(\mathsf{NULL}\) 的局部截面判据"与"多轴前缀提升可读性"本质上都属于对象层可断言性的组织问题。
为把这三条链压缩到同一语义框架中，本小节把对象层事实统一放回\emph{观察者索引的信息状态}上：全局层只保留推理规则，而对象层命题是否成立，则由局部信息状态上的 forcing 关系给出。

\begin{definition}[观察者索引的信息状态框架]\label{def:observer-indexed-information-frame}
称八元组
\[
\mathfrak F=(P,\le,\iota,\tau,\mathrm{Cov},\mathrm{Upd},D,V)
\]
为一个观察者索引的信息状态框架，若满足下列条件：
\begin{enumerate}
  \item \(P\) 为信息状态集合；每个 \(p\in P\) 表示一个局部可访问的信息条件，而非无索引的全局世界。
  \item \(\le\) 为 \(P\) 上的预序；若 \(q\le p\)，则称 \(q\) 是 \(p\) 的细化，表示 \(q\) 在 \(p\) 的基础上拥有更多记录、更细的区分能力或更强的分类结构。
  \item \(\iota:P\to I\) 为观察者索引映射，\(\tau:P\to T\) 为时间索引映射，并满足
  \[
  q\le p \Longrightarrow \iota(q)=\iota(p),\qquad \tau(q)\ge \tau(p).
  \]
  \item 对每个 \(p\in P\)，\(\mathrm{Cov}(p)\) 给出 \(p\) 的覆盖族；每个覆盖
  \[
  \{p_k\to p\}_{k\in K}
  \]
  满足 \(p_k\le p\)，并且覆盖在细化下稳定：若 \(q\le p\) 且 \(\{p_k\to p\}_{k\in K}\in\mathrm{Cov}(p)\)，则存在 \(\{q_\ell\to q\}_{\ell\in L}\in\mathrm{Cov}(q)\) 及映射 \(\lambda:L\to K\)，使得 \(q_\ell\le p_{\lambda(\ell)}\) 对一切 \(\ell\in L\) 成立。
  \item 对每个更新动作 \(\alpha\)，\(\mathrm{Upd}_\alpha(p)\subseteq P\) 为在 \(p\) 上执行 \(\alpha\) 后可到达的合法后继状态，并满足与细化的双向相容性：若 \(q\le p\)，则
  \[
  r\in \mathrm{Upd}_\alpha(q)\Rightarrow \exists s\in \mathrm{Upd}_\alpha(p)\ \text{使}\ r\le s,
  \]
  \[
  s\in \mathrm{Upd}_\alpha(p)\Rightarrow \exists r\in \mathrm{Upd}_\alpha(q)\ \text{使}\ r\le s.
  \]
  \item \(D\) 为对象域；对每个原子命题 \(A\)，赋值 \(V(A)\subseteq P\) 为下闭集，即
  \[
  p\in V(A),\ q\le p \Longrightarrow q\in V(A).
  \]
\end{enumerate}
\end{definition}

\begin{definition}[forcing 关系]\label{def:observer-indexed-forcing}
在框架 \(\mathfrak F\) 上，定义对象层可断言关系
\[
p\Vdash \varphi
\]
如下：
\begin{enumerate}
  \item 对原子命题 \(A\)，规定
  \[
  p\Vdash A \Longleftrightarrow p\in V(A).
  \]
  \item 对合取，规定
  \[
  p\Vdash \varphi\land\psi
  \Longleftrightarrow
  p\Vdash \varphi\ \text{且}\ p\Vdash \psi.
  \]
  \item 对蕴含，规定
  \[
  p\Vdash \varphi\to\psi
  \Longleftrightarrow
  \forall q\le p,\ \bigl(q\Vdash \varphi \Rightarrow q\Vdash \psi\bigr).
  \]
  \item 对否定，规定
  \[
  p\Vdash \neg\varphi
  \Longleftrightarrow
  \forall q\le p,\ q\nVdash \varphi.
  \]
  \item 对析取，规定 \(p\Vdash \varphi\lor\psi\) 当且仅当存在某个覆盖 \(\{p_k\to p\}_{k\in K}\in\mathrm{Cov}(p)\)，使得对每个 \(k\in K\)，均有
  \[
  p_k\Vdash \varphi\quad \text{或}\quad p_k\Vdash \psi.
  \]
  \item 对全称量词，规定
  \[
  p\Vdash \forall x\,\varphi(x)
  \Longleftrightarrow
  \forall q\le p,\ \forall d\in D,\ q\Vdash \varphi(d).
  \]
  \item 对存在量词，规定 \(p\Vdash \exists x\,\varphi(x)\) 当且仅当存在某个覆盖 \(\{p_k\to p\}_{k\in K}\in\mathrm{Cov}(p)\)，并可为每个 \(k\in K\) 选取 \(d_k\in D\)，使得
  \[
  p_k\Vdash \varphi(d_k).
  \]
  \item 对更新模态，规定
  \[
  p\Vdash [\alpha]\varphi
  \Longleftrightarrow
  \forall q\in \mathrm{Upd}_\alpha(p),\ q\Vdash \varphi,
  \]
  \[
  p\Vdash \langle\alpha\rangle\varphi
  \Longleftrightarrow
  \exists q\in \mathrm{Upd}_\alpha(p),\ q\Vdash \varphi.
  \]
\end{enumerate}
\end{definition}

\begin{definition}[已决定]\label{def:observer-indexed-decided}
称命题 \(\varphi\) 在状态 \(p\) 上\emph{已决定}，若
\[
p\Downarrow \varphi
\quad:\Longleftrightarrow\quad
p\Vdash \varphi\ \text{或}\ p\Vdash \neg\varphi.
\]
若既无 \(p\Vdash \varphi\)，又无 \(p\Vdash \neg\varphi\)，则称 \(\varphi\) 在 \(p\) 上未决定。
\end{definition}

\begin{proposition}[forcing 的单调性]\label{prop:observer-indexed-forcing-monotonicity}
若 \(q\le p\) 且 \(p\Vdash \varphi\)，则 \(q\Vdash \varphi\)。
\end{proposition}
\leanverified{paper\_recursive\_addressing\_observer\_indexed\_forcing\_monotonicity}

\begin{proof}
对公式复杂度作归纳。
原子情形由 \(V(A)\) 的下闭性直接得到；合取情形平凡。
蕴含、否定与全称量词情形都只需注意：定义中对一切 \(r\le p\) 的量化在 \(q\le p\) 下仍有效。
析取与存在量词情形由覆盖在细化下的稳定性推出：\(p\) 上的任一合法覆盖都可沿 \(q\le p\) 限制为 \(q\) 的合法覆盖。
对更新模态，\([\alpha]\) 的情形由后继状态与细化的前向相容性推出；\(\langle\alpha\rangle\) 的情形则由后继状态与细化的后向相容性推出。
故结论成立。
\end{proof}

\begin{corollary}[细化不推翻已建立的对象事实]\label{cor:observer-indexed-no-retraction}
若某命题在较粗状态 \(p\) 上已可断言，则在任何后续细化状态 \(q\le p\) 上仍可断言。
因此，信息细化的作用不是撤销对象层事实，而是缩小未决定域。
\end{corollary}
\leanverified{paper\_recursive\_addressing\_observer\_indexed\_no\_retraction}

\begin{proof}
这是命题 \ref{prop:observer-indexed-forcing-monotonicity} 的直接推论。
\end{proof}

\begin{proposition}[全局层只统一推理规则]\label{prop:observer-indexed-global-layer-logic-only}
设 \(\Gamma\vdash \varphi\) 是与定义 \ref{def:observer-indexed-forcing} 相容的演算层推理。
则对任意状态 \(p\in P\)，若
\[
p\Vdash \gamma
\qquad (\forall \gamma\in \Gamma),
\]
则
\[
p\Vdash \varphi.
\]
因此，全局层统一的是推理规则，而对象层事实始终由具体信息状态上的 forcing 给出。
\end{proposition}
\leanverified{paper\_recursive\_addressing\_observer\_indexed\_global\_layer\_logic\_only}

\begin{proof}
这是 forcing 语义对相应演算层的标准可靠性陈述；证明对推导长度作归纳即可。
证毕。
\end{proof}

\begin{definition}[地址对象与局部截面]\label{def:observer-indexed-address-local-sections}
沿用定义 \ref{def:prefix-site}--\ref{def:prefix-sheaf-readout} 的前缀站点与协议可接受读出层。
对每个信息状态 \(p\)，记
\[
A(p)\subseteq \mathcal{A}_{\mathrm{adm}}
\]
为在 \(p\) 上已对象化且协议允许的地址集合。
对每个 \(a\in A(p)\)，记
\[
\Gamma_p(a)\subseteq \mathscr{F}(a)
\]
为在状态 \(p\) 上、地址 \(a\) 处全部协议允许、类型意义下单值且可核验的局部截面之集合。
若 \(a\notin A(p)\)，则称该地址在 \(p\) 上尚未对象化；若 \(a\in A(p)\) 但 \(\Gamma_p(a)=\varnothing\)，则称该地址在 \(p\) 上尚无合法局部读出。
\end{definition}

\begin{proposition}[地址先于取值]\label{prop:observer-indexed-address-before-value}
对任意状态 \(p\) 与地址符号 \(a\)，若 \(a\notin A(p)\)，则一切形如"\(a\) 的值为 \(v\)"的对象层命题在 \(p\) 上都不能成立；
若 \(a\in A(p)\) 但 \(\Gamma_p(a)=\varnothing\)，则"\(a\) 在当前类型中可被单值读出"的命题在 \(p\) 上不能成立。
\end{proposition}
\leanverified{paper\_observer\_indexed\_address\_before\_value\_seeds}
\leanverified{paper\_observer\_indexed\_address\_before\_value\_package}
\leanverified{paper\_observer\_indexed\_address\_before\_value}

\begin{proof}
对象层取值命题要成立，前提是存在合法指称对象与相应的局部截面。
若地址对象尚未形成，则命题缺乏对象语义支撑；若地址已对象化但局部截面为空，则读出命题缺乏证书支撑。
因此两种情形下相关命题都不能在 \(p\) 上被强制。
\end{proof}

\begin{corollary}[非 \texorpdfstring{$\mathsf{NULL}$}{NULL} 的 forcing 判据]\label{cor:observer-indexed-nonnull-criterion}
在状态 \(p\) 上，地址 \(a\) 的读出为非 \(\mathsf{NULL}\) 当且仅当
\[
a\in A(p)
\qquad\text{且}\qquad
\Gamma_p(a)\neq \varnothing.
\]
因此，非 \(\mathsf{NULL}\) 不是一个值域元素，而是局部截面的存在性条件。
\end{corollary}
\leanverified{paper\_recursive\_addressing\_observer\_indexed\_nonnull\_criterion}

\begin{proof}
由命题 \ref{prop:observer-indexed-address-before-value} 立即得到。
\end{proof}

\begin{definition}[\texorpdfstring{$\mathsf{NULL}$}{NULL} 的状态化三分解]\label{def:observer-indexed-null-trichotomy}
在状态 \(p\) 上，地址 \(a\) 的 \(\mathsf{NULL}\) 分为三类：
\begin{enumerate}
  \item \textbf{语义空值：} \(a\notin A(p)\)。此时对象层尚未形成合法指称。
  \item \textbf{协议空值：} \(a\in A(p)\) 但由于当前可见域、类型约束或协议纪律的限制，\(\Gamma_p(a)=\varnothing\)。此时有地址而无合法读出。
  \item \textbf{碰撞空值：} 地址与协议均合法，但底层读出仍存在不可消解的多对一压缩；在未补足足够残差轴、未提升到单射层级之前，\(\Gamma_p(a)=\varnothing\)。
\end{enumerate}
\end{definition}

\begin{proposition}[\texorpdfstring{$\mathsf{NULL}$}{NULL} 的结构性]\label{prop:observer-indexed-null-structural}
\(\mathsf{NULL}\) 不是值域中的一个普通元素，而是对象层 forcing 失败的结构性形式。
\leanverified{paper\_recursive\_addressing\_observer\_indexed\_null\_structural}
\end{proposition}

\begin{proof}
由推论 \ref{cor:observer-indexed-nonnull-criterion}，非 \(\mathsf{NULL}\) 的充要条件是局部截面的非空存在。
故 \(\mathsf{NULL}\) 的本质不是某个特殊值被读出，而是相关读出命题未能在当前状态上被强制成立。
\end{proof}

\begin{definition}[多轴前缀与公共细化]\label{def:observer-indexed-common-refinement}
设 \(u\) 与 \(v\) 是两个协议兼容的前缀轴，并记 \(p_u,p_v\in P\) 为分别只包含单轴 \(u,v\) 信息的状态。
若存在状态 \(p_{u\otimes v}\in P\) 使
\[
p_{u\otimes v}\le p_u,
\qquad
p_{u\otimes v}\le p_v,
\]
则称 \(p_{u\otimes v}\) 为它们的联合前缀状态或公共细化。
\end{definition}

\begin{proposition}[多轴前缀提升可断言性]\label{prop:observer-indexed-common-refinement-decidability}
\leanverified{paper\_recursive\_addressing\_observer\_indexed\_common\_refinement\_decidability}
在定义 \ref{def:observer-indexed-common-refinement} 的前提下，联合前缀不会降低任何已建立的对象层事实。
并且，凡依赖于额外残差轴才能消解的碰撞命题，都可能在公共细化状态上首次变为已决定。
\end{proposition}

\begin{proof}
第一部分是命题 \ref{prop:observer-indexed-forcing-monotonicity} 的直接结果。
第二部分则来自公共细化对不可区分候选读出的进一步切分：原先在单轴状态上混合在一起的候选项，经附加轴区分后，可能首次形成单值且可核验的局部截面，从而把未决定命题推进到已决定状态。
\end{proof}

\begin{definition}[延迟分类]\label{def:observer-indexed-delayed-classification}
设 \(s\le \tau(p)\)，且 \(A_s\) 为一个关于较早时刻 \(s\) 的命题。
若在状态 \(p\) 上
\[
p\nVdash A_s,
\qquad
p\nVdash \neg A_s,
\]
但存在某个更新动作 \(\alpha\) 与某个后继状态 \(q\in \mathrm{Upd}_\alpha(p)\)，使得
\[
q\Downarrow A_s,
\]
则称 \(A_s\) 在 \(p\to q\) 的更新过程中发生了延迟分类。
\end{definition}

\begin{proposition}[延迟分类不改写过去]\label{prop:observer-indexed-delayed-classification-not-retrocausal}
在定义 \ref{def:observer-indexed-delayed-classification} 的情形下，更新所改变的是关于时刻 \(s\) 之命题的当前可断言性，而不是时刻 \(s\) 的对象层内容本身。
\end{proposition}

\begin{proof}
命题 \(A_s\) 的时间索引 \(s\) 在更新前后保持不变；变化的只是信息状态由 \(p\) 细化为某个可达后继 \(q\)。
因此，延迟分类表达的是"关于过去对象的命题在更强信息状态上首次被定解"，而非"过去被回写为另一个对象层事实"。
\end{proof}
\leanverified{paper\_recursive\_addressing\_observer\_indexed\_delayed\_classification\_not\_retrocausal}

\begin{corollary}[观察更新的对象层意义]\label{cor:observer-indexed-update-decides-not-create}
观察、再分类、增加残差轴、拼接证书对象等操作，在对象层上的统一作用，是把某些命题从未决定推进到已决定，而不是凭空创造事实。
\leanverified{paper\_recursive\_addressing\_observer\_indexed\_update\_decides\_not\_create}
\end{corollary}

\begin{proposition}[递归柱事件的有限构造证书与 incidence 账本]\label{prop:distill-euclid-elements-recursive-cylinder-construction-certificate}
\textbf{Dependency status: construction.}
沿用命题 \ref{prop:recursive-addressing-refinement} 的记号。给定
\[
a=(a_0,\ldots,a_{\ell-1})\in\bigl(\{0,1\}^{r_L}\bigr)^\ell,
\qquad
a_j=(a_{j,1},\ldots,a_{j,r_L}),
\]
定义有符号原子
\[
E_{j,i}(a):=
\begin{cases}
T^{-(t+j)}C_i^{(L)},& a_{j,i}=1,\\
X\setminus T^{-(t+j)}C_i^{(L)},& a_{j,i}=0.
\end{cases}
\]
则有限列表
\[
\mathfrak C(L,a,t):=\bigl(E_{j,i}(a)\bigr)_{0\le j<\ell,\ 1\le i\le r_L}
\]
是 \(C^{(L+1)}_{a,t}\) 的构造证书，并且其 incidence 账本满足：
\begin{enumerate}
  \item 精确构造式
  \[
  C^{(L+1)}_{a,t}
  =\bigcap_{j=0}^{\ell-1}\bigcap_{i=1}^{r_L}E_{j,i}(a);
  \]
  \item 若 \(b=(b_0,\ldots,b_{\ell'-1})\) 满足 \(\ell'\ge\ell\) 且 \(b_j=a_j\) 对 \(0\le j<\ell\) 成立，则
  \[
  C^{(L+1)}_{b,t}\subseteq C^{(L+1)}_{a,t};
  \]
  \item 若 \(a,a'\in(\{0,1\}^{r_L})^\ell\) 且存在 \((j,i)\) 使 \(a_{j,i}\ne a'_{j,i}\)，则
  \[
  C^{(L+1)}_{a,t}\cap C^{(L+1)}_{a',t}=\varnothing.
  \]
\end{enumerate}
因此，递归地址化中的派生概念先由有限旧层原子构造，再由包含或相交关系进行谓词化；任何 incidence 断言均可回溯到同一张有限 Boolean 账本。
\end{proposition}

% --- Distillation writeback ---
\begin{theorem}[地址歧义的条件熵降与无界寄存器泄漏]\label{distill:bourgain-entropy_increment_closure-address-ambiguity-entropy-leak}
固定状态 \(p\) 与已对象化地址 \(a\in A(p)\)。设
\[
E\subseteq \Gamma_p(a)
\]
为一族有限候选局部截面，\(\mathfrak m\) 为 \(E\) 上满支撑概率测度，\(\rho:E\to Y\) 为当前协议已经登记的有限可见读出。用自然对数定义 Shannon 熵，并记
\[
\mathsf H_\rho(E):=H_{\mathfrak m}(E\mid \rho)
=
\sum_{y\in\rho(E)}\mathfrak m(\rho^{-1}(y))H(\mathfrak m_y),
\]
其中 \(\mathfrak m_y\) 是 \(\rho^{-1}(y)\) 上的条件测度。称有限映射 \(\chi:E\to Z\) 为相对于 \(\rho\) 的授权细化判别器，若其图可由某个共同细化状态 \(q\le p\) 登记，且不改变地址对象 \(a\)。令
\[
\rho^\chi:E\longrightarrow Y\times Z,
\qquad
\rho^\chi(e):=(\rho(e),\chi(e)).
\]
则：
\begin{enumerate}
  \item 条件熵满足精确降公式
  \[
  \mathsf H_{\rho^\chi}(E)
  =\mathsf H_\rho(E)-H_{\mathfrak m}(\chi\mid \rho).
  \]
  特别地，\(\mathsf H_{\rho^\chi}(E)<\mathsf H_\rho(E)\) 当且仅当存在正 \(\mathfrak m\)-质量的 \(\rho\)-纤维使 \(\chi\) 在其上非常值。
  \item 若 \(\rho^\chi\) 在 \(\mathfrak m\) 的支撑上为终端读出，即每个正质量 \(\rho^\chi\)-纤维均为单点，则地址 \(a\) 在对应细化状态 \(q\) 上由 \(E\) 给出单值且非 \(\mathsf{NULL}\) 的局部读出证书。
  \item 设 \((E_N,\mathfrak m_N,\rho_N)_{N\ge0}\) 是同一地址 \(a\) 的一列有限歧义审计，并假设每个授权细化判别器都在每条正 \(\mathfrak m_N\)-质量的 \(\rho_N\)-纤维上为常值。置
  \[
  M_N:=\max_{y\in\rho_N(E_N)}\bigl|E_N\cap\rho_N^{-1}(y)\bigr|.
  \]
  若 \(\sup_N M_N=\infty\)，则满足如下量化泄漏性质：对每个 \(B\in\NN\)，存在 \(N\) 与一条正质量纤维 \(F\subseteq E_N\)，使得任意 \(|R|\le B\) 的有限寄存器读出 \(r:F\to R\) 都不能把 \(F\) 终端分离。
\end{enumerate}
\end{theorem}
\begin{proof}
因为 \(\chi\) 是 \(E\) 的函数，条件熵链式法则给出
\[
H_{\mathfrak m}(E\mid \rho)
=H_{\mathfrak m}(\chi,E\mid \rho)
=H_{\mathfrak m}(\chi\mid \rho)+H_{\mathfrak m}(E\mid \rho,\chi).
\]
而 \(H_{\mathfrak m}(E\mid \rho,\chi)=\mathsf H_{\rho^\chi}(E)\)，于是得到第一条公式。有限概率空间上，条件熵 \(H_{\mathfrak m}(\chi\mid \rho)\) 为零当且仅当在每条正质量 \(\rho\)-纤维上 \(\chi\) 的条件分布为 Dirac 测度，也即 \(\chi\) 为常值；严格性判据随之成立。

若 \(\rho^\chi\) 的每条正质量纤维均为单点，则每个可见读出值至多对应一个候选局部截面。由于 \(E\subseteq\Gamma_p(a)\)，且 \(q\le p\) 只是在同一地址对象上登记额外判别信息，已存在的局部截面证书不因细化而被撤销。由推论 \ref{cor:observer-indexed-nonnull-criterion}，地址已对象化且存在唯一可核验局部截面时给出非 \(\mathsf{NULL}\) 读出，第二条成立。

对第三条，任取 \(B\in\NN\)。由 \(\sup_N M_N=\infty\)，可取 \(N\) 与 \(y\in\rho_N(E_N)\) 使
\[
F:=E_N\cap\rho_N^{-1}(y),
\qquad |F|>B.
\]
若某个 \(|R|\le B\) 的寄存器读出 \(r:F\to R\) 终端分离 \(F\)，则 \(r\) 必须在 \(F\) 上单射；否则两个不同候选截面具有相同寄存器值而仍未被区分。单射性给出 \(|R|\ge |F|>B\)，与 \(|R|\le B\) 矛盾。假设中授权细化判别器均在正质量歧义纤维上为常值，故这类纤维不能通过内部细化产生第一条中的严格熵降。量化泄漏性质得证。
\end{proof}

% --- End distillation writeback ---

\begin{proof}
第一条只是把条件
\[
\mathbf b^{(L)}_{t+j}(x)=a_j
\qquad(0\le j<\ell)
\]
按坐标展开。对固定 \((j,i)\)，等式的第 \(i\) 个坐标为 \(1\) 时要求 \(x\in T^{-(t+j)}C_i^{(L)}\)，为 \(0\) 时要求 \(x\notin T^{-(t+j)}C_i^{(L)}\)。对全部 \((j,i)\) 取有限交即得构造式。

若 \(b\) 延长 \(a\)，则 \(C^{(L+1)}_{b,t}\) 的定义包含 \(C^{(L+1)}_{a,t}\) 的全部原子条件，并额外加入更深时刻的条件，故有包含关系。若 \(a\) 与 \(a'\) 在某个坐标 \((j,i)\) 不同，则一个柱事件要求 \(x\in T^{-(t+j)}C_i^{(L)}\)，另一个要求 \(x\notin T^{-(t+j)}C_i^{(L)}\)，二者不相交。证毕。
\end{proof}

\begin{proof}
由命题 \ref{prop:observer-indexed-common-refinement-decidability} 与命题 \ref{prop:observer-indexed-delayed-classification-not-retrocausal} 联立即得。
\end{proof}

\paragraph{类型化读出偏函数与对象层值}
前述 forcing 框架还可进一步压缩为一个状态化的类型读出偏函数。
它的作用是把“地址是否已对象化”“局部截面是否存在”与“值类是否已在当前类型制度中唯一确定”三件事统一到同一个对象层接口上。

\begin{definition}[类型化读出偏函数]\label{def:observer-indexed-typed-readout}
对任意信息状态 \(p\) 与类型化地址 \(a\)，定义偏函数 \(\mathrm{Read}_p(a)\) 如下：
\begin{enumerate}
  \item 若 \(a\notin A(p)\)，则 \(\mathrm{Read}_p(a)\) 在对象层不定义；
  \item 若 \(a\in A(p)\) 但 \(\Gamma_p(a)=\varnothing\)，则记
  \[
  \mathrm{Read}_p(a)=\mathsf{NULL};
  \]
  \item 若 \(a\in A(p)\) 且 \(\Gamma_p(a)\) 在当前类型制度下确定唯一可核验值类，则记该值类为 \(\mathrm{Read}_p(a)\)。
\end{enumerate}
当可见域进一步固定为 unitary slice 时，记相应特化为
\[
\mathrm{Read}_{\mathrm{US},p}(a).
\]
其跨状态的全局编译将在后续逆极限口径下组织为偏函数 \(\mathrm{Read}_{\mathrm{US}}\)，见定理 \ref{thm:recursive-addressing-readout-separatedness-null}。
\end{definition}

\begin{proposition}[\texorpdfstring{$\mathrm{Read}_{\mathrm{US}}$}{Read\_US} 的 forcing 判据]\label{prop:observer-indexed-readout-forcing-criterion}
\leanverified{paper\_recursive\_addressing\_observer\_indexed\_readout\_forcing\_criterion}
在当前类型纪律下，对任意状态 \(p\) 与类型化地址 \(a\)，有
\[
\mathrm{Read}_{\mathrm{US},p}(a)\neq \mathsf{NULL}
\]
当且仅当当前状态上存在协议允许、类型单值且可核验的局部截面。
并且若 \(q\le p\) 且 \(\mathrm{Read}_{\mathrm{US},p}(a)\neq \mathsf{NULL}\)，则
\[
\mathrm{Read}_{\mathrm{US},q}(a)\neq \mathsf{NULL},
\]
且其对象层值类与 \(\mathrm{Read}_{\mathrm{US},p}(a)\) 一致。
\end{proposition}

\begin{proof}
前半句由定义 \ref{def:observer-indexed-typed-readout} 直接给出。
对后半句，由命题 \ref{prop:observer-indexed-forcing-monotonicity}，一旦某个读出命题在较粗状态上已可断言，则在任何细化状态上仍可断言。
而定义 \ref{def:observer-indexed-typed-readout} 的第三条已把合法局部截面压缩为唯一可核验值类，因此细化只会增加证书支撑，不会改变对象层值类本身。
故结论成立。
\end{proof}

\begin{corollary}[显式提升原则的 forcing 解释]\label{cor:observer-indexed-explicit-lifting}
\leanverified{paper\_recursive\_addressing\_observer\_indexed\_explicit\_lifting}
若某个复合表达式同时调用来自不同精度索引、不同可见域或不同证书层级的读出，而系统并未提供一个共同细化状态 \(q\le p\)，使这些读出能在同一类型化域中同时解释，则该表达式在状态 \(p\) 上不定义。
\end{corollary}

\begin{proof}
复合表达式的对象层解释要求其各子项在同一可比较域中具有共同语义。
若不存在共同细化把这些子项提升到同一类型层，则不存在统一的对象层评价函数，因而该表达式不能在 \(p\) 上被强制成立。
这正是“混合精度的操作在类型层不定义”的 forcing 口径。
\end{proof}

\paragraph{值保持重写的对象层语义}
一旦地址对象已经形成、局部截面已经存在并且对象层值类已被确定，后续的稳定算术便不再处理“是否存在对象”，而只处理“如何在同一对象层值类内选取规范表示”。
相应地，值保持重写必须在对象层而非实现层被表述。

\begin{definition}[值保持重写的对象层语义]\label{def:observer-indexed-value-preserving-rewrite}
设 \(t\Rightarrow t'\) 是后续稳定算术中的一条重写规则。
称该规则在对象层上\emph{值保持}，若对任意信息状态 \(p\) 与任意使两侧都可解释的地址--值环境 \(\eta\)，均有
\[
p\Vdash \operatorname{Val}_{p,\eta}(t)=\operatorname{Val}_{p,\eta}(t').
\]
其中 \(\operatorname{Val}_{p,\eta}(t)\) 表示项 \(t\) 在状态 \(p\) 与环境 \(\eta\) 下的对象层读出类。
\end{definition}

\begin{proposition}[值保持重写不创造对象层值]\label{prop:observer-indexed-value-preserving-no-creation}
若一条重写规则在定义 \ref{def:observer-indexed-value-preserving-rewrite} 的意义下值保持，则它只能在已有对象层读出上改变表示，而不能从一个未定义项中创造出新的对象层值。
\leanverified{paper\_recursive\_addressing\_observer\_indexed\_value\_preserving\_no\_creation}
\end{proposition}

\begin{proof}
若 \(t\) 在状态 \(p\) 上未定义，则 \(\operatorname{Val}_{p,\eta}(t)\) 不存在。
此时若 \(t'\) 在 \(p\) 上已定义，则
\[
p\Vdash \operatorname{Val}_{p,\eta}(t)=\operatorname{Val}_{p,\eta}(t')
\]
无从成立，因此该规则不可能是对象层值保持的。
反之，若两侧都已定义，则值保持只说明二者属于同一对象层值类，因而改变的是表示、规范形或证书接口，而不是值本身。
\end{proof}

\begin{corollary}[规范化不创造事实]\label{cor:observer-indexed-normalization-no-fact-creation}
后续的规范化、重写与横截选择，其对象层作用只是把一个已定义值的不同表示压到统一规范形；它们不能替代地址对象化、证书建立或前缀细化，更不能越过 \(\mathsf{NULL}\) 直接生成数值事实。
\leanverified{paper_recursive_addressing_observer_indexed_normalization_no_fact_creation_seeds}
\end{corollary}

\begin{proof}
这是命题 \ref{prop:observer-indexed-value-preserving-no-creation} 的直接推论。
\end{proof}

\begin{proposition}[有限分辨率算术的对象层资格]\label{prop:observer-indexed-finite-resolution-object-eligibility}
\leanverified{paper\_recursive\_addressing\_observer\_indexed\_finite\_resolution\_object\_eligibility}
在分辨率 \(m\) 上，稳定算术只对已经通过对象化与非 \(\mathsf{NULL}\) 判据的读出类生效。
更具体地说，若
\[
\omega\equiv_m\omega'
\Longleftrightarrow
\Fold_m(\omega)=\Fold_m(\omega')
\]
给出实现层的折叠同余，则对象层算术并不直接作用于原始微态 \(\omega\)，而作用于由当前状态所支撑的读出类，即实现层等价类经地址化与类型过滤之后的对象。
\end{proposition}

\begin{proof}
折叠同余只说明不同微态在当前分辨率下不可区分；它本身既不构成对象层地址，也不自动给出可核验读出。
只有当相应等价类在当前信息状态上被地址化、通过协议并获得局部截面时，它才上升为对象层对象。
故稳定算术处理的不是原始微态，而是 forcing 已经承认的读出类。
\end{proof}

\begin{corollary}[第 \texorpdfstring{\ref{sec:emergent-arithmetic}}{7} 节的读法]\label{cor:observer-indexed-reading-of-emergent-arithmetic}
\leanverified{paper\_recursive\_addressing\_observer\_indexed\_reading\_of\_emergent\_arithmetic}
第 \ref{sec:emergent-arithmetic} 节中的值保持重写、规范横截面与分辨率模结构，应理解为对那些已经在本节意义下完成对象化、并已获得非 \(\mathsf{NULL}\) 读出的地址--值对，构造其规范表示与值保持演算。
更具体地说：
\begin{enumerate}
  \item 第 \ref{sec:emergent-arithmetic} 节并不是在无地址对象上直接施行算术；
  \item 其中的“值保持”并不是实现层字符串恒等，而是对象层读出类恒等；
  \item 其中的“模结构”也不是任意商，而是由分辨率折叠与可核验读出共同决定的对象层商结构。
\end{enumerate}
\leanverified{paper\_recursive\_addressing\_observer\_indexed\_reading\_of\_emergent\_arithmetic}
\end{corollary}

% --- Distillation writeback ---
\begin{proposition}[对象层等式的有限 forcing 账本正规形]\label{distill:euclid_elements-finite_construction_certificate_families-forcing-equality-ledger-normal-form}
\textbf{日期时间：2026-04-23T11:58:48+08:00。依赖状态：rigidity.}
固定信息状态 \(p\)。考虑由下列 primitive 语法生成的对象层等式证明：
\begin{enumerate}
  \item 已在某个共同细化状态 \(q\le p\) 上非 \(\mathsf{NULL}\) 的读出 \(\mathrm{Read}_{\mathrm{US},q}(a)\)；
  \item 定义 \ref{def:observer-indexed-value-preserving-rewrite} 意义下已登记的值保持重写；
  \item common-notion 规则：自反、对称、传递以及对已定义对象层值的替换。
\end{enumerate}
称一份有限 equality ledger 为一棵有限有向无环证明树，其叶节点只允许是同一非 \(\mathsf{NULL}\) 读出的对象层同一性或已登记值保持重写，其内节点只允许使用上述 common-notion 规则。则在该语法内：
\begin{enumerate}
  \item 若参与复合表达式的各读出没有共同细化状态使其落在同一类型化域中，则该表达式在状态 \(p\) 上不定义；
  \item 任一有限 equality ledger 的根结论
  \[
  t=t'
  \]
  都满足
  \[
  p\Vdash \operatorname{Val}_{p,\eta}(t)=\operatorname{Val}_{p,\eta}(t')
  \]
  对所有使两侧已定义的地址--值环境 \(\eta\) 成立；
  \item 反过来，凡是在上述 primitive 语法中推出的对象层等式，都可整理为一份有限 equality ledger。因此，在该语法内，等式不能由图示相似性或规范化步骤凭空产生；它必须回指到某个非 \(\mathsf{NULL}\) 读出、某条值保持重写或一条 common-notion 规则。
\end{enumerate}
\end{proposition}

% --- Distillation writeback ---
\begin{proposition}[随机分支读出的确定性冻结与提升门控]\label{prop:distill-bourgain-forcing-random-certificate-extraction}
\textbf{依赖状态：conditional audit / conditional rigidity.}
固定信息状态 \(p\)，有限指标集 \(I\)，以及对每个 \(i\in I\) 的细化状态 \(q_i\le p\)。令 \(\mathcal E\) 表示一个对象层等式断言 \(t=t'\)。假设给定有限审计数据
\[
(\lambda_i,L_i,\sigma_i,\pi_i)_{i\in I},
\]
其中 \(\lambda_i\in\QQ_{\ge0}\)、\(\sum_i\lambda_i=1\)，\(L_i\) 是一份候选有限 equality ledger 或空符号，\(\sigma_i,\pi_i\in\{0,1\}\)。这些数据的语义规定如下：若 \(\sigma_i=1\)，则 \(L_i\) 是在状态 \(q_i\) 上证明 \(\mathcal E\) 的有限 equality ledger，因而对所有使两侧在 \(q_i\) 上已定义的地址--值环境 \(\eta\) 有
\[
q_i\Vdash \operatorname{Val}_{q_i,\eta}(t)=\operatorname{Val}_{q_i,\eta}(t').
\]
若 \(\pi_i=1\)，则验证器还显式登记了提升规则
\[
(q_i,L_i,\mathcal E)\leadsto_p \mathcal E,
\]
其含义正是
\[
p\Vdash \operatorname{Val}_{p,\eta}(t)=\operatorname{Val}_{p,\eta}(t')
\]
对所有使两侧在 \(p\) 上已定义的环境 \(\eta\) 成立。定义该有限随机分支族的可提升成功质量为
\[
M_p(I):=\sum_{i\in I}\lambda_i\sigma_i\pi_i.
\]
则：
\begin{enumerate}
  \item 若 \(M_p(I)>0\)，则存在 \(i_\ast\in I\) 使 \(\lambda_{i_\ast}>0\)、\(\sigma_{i_\ast}=1\)、\(\pi_{i_\ast}=1\)。冻结三元组 \((q_{i_\ast},L_{i_\ast},\mathcal E)\) 后得到状态 \(p\) 上的一份确定性对象层等式证书；
  \item 若 \(\sum_i\lambda_i\sigma_i>0\)，则至少存在某个细化状态 \(q_i\) 上的确定性等式 ledger；但除非相应 \(π_i=1\)，该 ledger 只停留在 \(q_i\) 层，不能自动升级为 \(p\) 上的 forcing 结论；
  \item 若 \(M_p(I)=0\)，则只能推出当前抽样分支族中没有同时满足成功与提升登记的分支。它不推出 \(p\) 上 \(\mathcal E\) 为假，也不排除存在未列入 \(I\) 的其它确定性证明账本。
\end{enumerate}
\end{proposition}
\begin{proof}
第一条由有限和的正性直接给出：若每个正权分支都满足 \(\sigma_i\pi_i=0\)，则 \(M_p(I)=0\)，与假设矛盾。因此可取 \(i_\ast\) 使 \(\lambda_{i_\ast}>0\) 且 \(\sigma_{i_\ast}=\pi_{i_\ast}=1\)。由 \(\sigma_{i_\ast}=1\)，\(L_{i_\ast}\) 是 \(q_{i_\ast}\) 上的有限 equality ledger；再由 \(\pi_{i_\ast}=1\) 的定义，登记的提升规则把该 ledger 编译为状态 \(p\) 上的对象层等式证书。

第二条同理，\(\sum_i\lambda_i\sigma_i>0\) 给出一个正权成功分支，但该结论只使用 \(\sigma_i\) 的语义，因此只得到 \(q_i\) 层 forcing。forcing 的通常细化单调性并不提供从 \(q_i\le p\) 回到 \(p\) 的反向提升；反向提升恰由 \(\pi_i=1\) 单独登记。

第三条只是对有限审计域的否定陈述。\(M_p(I)=0\) 说明在已列出的有限分支中没有正权分支同时通过成功位与提升位；该命题没有量化所有可能 ledger，故不能推出对象层断言在 \(p\) 上不可证或为假。证毕。
\end{proof}

% --- End distillation writeback ---

\begin{proof}
第一条正是推论 \ref{cor:observer-indexed-explicit-lifting}：复合表达式要求所有子项在同一可比较域中解释；若不存在共同细化状态，则没有统一对象层评价函数，故表达式不定义。

第二条对 ledger 的高度作归纳。高度为零时，叶节点分两类。若叶节点是同一非 \(\mathsf{NULL}\) 读出的对象层同一性，则由定义 \ref{def:observer-indexed-typed-readout} 的第三条，读出类已经唯一确定，故等式成立。若叶节点是已登记的值保持重写，则结论正是定义 \ref{def:observer-indexed-value-preserving-rewrite}。归纳步中，自反、对称与传递保持对象层等式；替换规则只在已定义对象层值上使用，因而由等式的相容性保持强制成立。故 ledger 根结论被 \(p\) 强制。

第三条由语法本身的有限生成性给出。任一推导只有有限多步；把其原始读出同一性与值保持重写作为叶节点，把每次 common-notion 应用作为内节点，即得到有限有向无环证明树。命题 \ref{prop:observer-indexed-value-preserving-no-creation} 进一步保证值保持重写不能从未定义项创造对象层值；因此任何合法等式推导都不可能越过 \(\mathsf{NULL}\) 或未构造地址而生成事实。证毕。
\end{proof}

% --- End distillation writeback ---

\begin{proof}
由推论 \ref{cor:observer-indexed-nonnull-criterion}，对象层算术的合法输入首先必须是非 \(\mathsf{NULL}\) 的地址--值对；
由命题 \ref{prop:observer-indexed-value-preserving-no-creation}，值保持重写只能在已定义对象上改变表示；
再由命题 \ref{prop:observer-indexed-finite-resolution-object-eligibility}，有限分辨率下真正进入算术层的对象，是折叠同余经对象化与类型过滤之后的读出类。
三者合并即得结论。
\end{proof}

\paragraph{本节收束}
至此，第 \ref{sec:recursive-addressing} 节可压缩为一条统一逻辑链：
先由地址生成把对象带入可指称域；
再由前缀站点、局部截面与证书对象决定何者可读；
再由 forcing 语义区分成立、失败与未决定；
再由多轴前缀与显式提升把未决定命题推进到可断言；
最后，只有在这一切都完成之后，稳定算术才获得其合法适用域。
因此，第 \ref{sec:emergent-arithmetic} 节并不是从零开始引入一个外加层，而是本节对象层语义的代数化延续：
本节回答“何时可以谈值”，而第 \ref{sec:emergent-arithmetic} 节回答“在已经可以谈值之后，如何进行值保持的规范化与算术化”。


